\cotefile{161-2}
\scanpage{66}
Existence de $\varprojlim$ ou $\varinjlim$ dans $T(C)\inserted{\simeq\underline{\mathrm{Hom}}_{\Delta}(R,C)}$, pour $T$ \unsure{2-représentable|2-representable}.

Soit $\delta$ un type de diagramme, tel que

a) Les $\varprojlim$ (disons) de type $\delta$ existent dans $\mathbf C$

b) les $\varprojlim$ de type $\delta$ commutent, \deleted{pour tt $C'\in\mathrm{Ob}\,\Delta$,} \inserted{dans $C$} aux $\unsure{\lim|\varinjlim}$ de type $d$ (i.e. aux $\varinjlim$ de type $d''$)

Alors les $\varprojlim$ de type $\delta$ existent dans $T(C)$, et se calculent dans $\underline{\mathrm{Hom}}_{\Delta}(R,C)$ ``\uline{argument} par \uline{argument}'' (en d'autres termes, les foncteurs $r_{C}:T(C)\longrightarrow C$ associés aux \inserted{$T$-}constructions commutent aux $\varprojlim$ de type $\delta$)

De plus, pour un hom de théories \unsure{repr.} $T\longrightarrow T'$, $T(C)\longrightarrow T'(C)$ commute aux $\varprojlim$ de type $d$. Id pour $C\longrightarrow C'$, $C'$ satisfaisant les $\unsure{\text{mêmes}|\leftrightarrows}$ hyp. que $C$, $T(C)\longrightarrow T(C')$ commute aux $\varprojlim$ de type $d$.

On a, pour $R,\unsure{\mathscr C|\mathbf C}\in\mathrm{Ob}\,\Delta$, \deleted{\unsure{$\mathrm{Ob}\,\mathcal C$}} et $\unsure{\mathbf J|\mathbf I}\in\mathrm{Ob}\,\mathrm{Cat}$
\[
\underline{\mathrm{Hom}}_{\mathbf{Cat}}(\mathbf J,\mathscr C)\in\Delta
\qquad\text{(au moins si $\mathbf J$ petite)}
\]
\[
\underline{\mathrm{Hom}}_{\Delta}(R,\underline{\mathrm{Hom}}_{\mathrm{Cat}}(\mathbf J,\mathscr C))
\simeq\underline{\mathrm{Hom}}_{\mathrm{Cat}}(\mathbf J,\underline{\mathrm{Hom}}_{\Delta}(R,\mathscr C))
\]
i.e.
\[
\boxed{T_R(\underline{\mathrm{Hom}}_{\mathrm{Cat}}(\mathbf J,\mathscr C))\simeq
\underline{\mathrm{Hom}}_{\mathrm{Cat}}(\mathbf J,T_R(C))}
\]

Faisons p. ex. $\mathscr C=(\mathrm{Ens})$ (supposant $(\mathrm{Ens})\in\mathrm{Ob}\,\Delta$ !), $\mathbf J$ remplacé par $\mathbf J^{\circ}$
\[
T_{\unsure{\mathscr R|R}}(\widehat{\mathbf J})\overset{\sim}{\simeq}\underline{\mathrm{Hom}}(J^{\circ},T_R(\mathrm{Ens}))
\]

Si $C\longrightarrow C'$ dans $\Delta$ est pl. fidèle (resp. fid) alors $T(C)\longrightarrow T(C')$ l'est aussi (pour $T$ 2-repr.). \deleted{et} Si $T\longrightarrow T'$ est associé à $R'\longrightarrow R^{\unsure{\prime\prime|\prime}}$ ``\inserted{$\Delta$-}générateur'' (tel que, donc, tte sous-catégorie pleine $\in\Delta$ de $R'$ par \unsure{lequel|laquel|laquelle} ça se factorise soit \deleted{\unsure{dans}} $\simeq R$] alors le diagramme
\diagram{\begin{tikzcd}[ampersand replacement=\&]
T(C) \arrow[r,hook] \arrow[d] \& T(C') \arrow[d] \\
T'(C) \arrow[r,hook] \& T'(C')
\end{tikzcd}}
est 2-cartésien
% archive sheet number: 65
\scanpage{67}
[Cas particulier \unsure{où|si} $T'=\tau^{I}$ ....]
% archive sheet number: 66
\scanpage{68}
\margin{\textcircled{\unsure{2|*|4}}}
% slanted
\margin{\uline{Cas}}
$d''=\varnothing$ : structures alg. définies seulement par $\varprojlim$ : Faits particuliers :

1°) $\exists$ $\varprojlim$ quelc. dans les $T(C)$, les $T(C)\longrightarrow T(C')$ y commutent, de \unsure{même|$\simeq$} les $T(C)\longrightarrow T'(C)$ ($T,T'$ 2-repr.), ainsi que les $T(C)\xrightarrow{r}C$ (foncteurs constructions).

2°) $C\longrightarrow\widehat C$ est un $\Delta$-foncteur, \deleted{\unsure{l'exactitude de|l'exactit. de}} (si $(\mathrm{Ens})\in\mathrm{Ob}\,\Delta$) donc on trouve que (si $T$ 2-rep par $R$)
\[
T_{\deleted{\unsure{R}}}(C)\longrightarrow T_{\deleted{R}}(\widehat C)
\simeq\underline{\mathrm{Hom}}_{\mathrm{Cat}}(C^{\circ},\underset{\substack{\Vert\\\inserted{S}}}{T_{\deleted{R}}(\mathrm{Ens})})
\]
est pl. fidèle, et identifie $T_{\deleted{R}}(C)$ à la sous-catégorie str. pleine de $\underline{\mathrm{Hom}}_{\mathrm{Cat}}(C^{\circ},S)$ formée des $\varphi:C^{\circ}\longrightarrow S$ tels que, pour tte \uline{construction} $r\in R$, \inserted{d'où $\widetilde r:S\longrightarrow(\mathrm{Ens})$,} $\widetilde r\circ\varphi:C^{\circ}\longrightarrow(\mathrm{Ens})$ \unsure{sont|soit} repr. Et il suffit d'exiger cette condition pour les $r_i$ qui $\Delta$-\uline{engendrent} (i.e. on a un diagramme 2-cartésien
\diagram{\begin{tikzcd}[ampersand replacement=\&,column sep=large]
T(C) \arrow[r] \arrow[d] \& \underline{\mathrm{Hom}}_{\mathrm{Cat}}(C^{\circ},T(\mathrm{Ens})) \arrow[d,"{\varphi\mapsto(\deleted{\unsure{\varphi}}\widetilde r_i\circ\varphi)}"] \\
C^{I} \arrow[r] \& \underline{\mathrm{Hom}}_{\mathrm{Cat}}(C^{\circ},(\mathrm{Ens}))^{I}
\end{tikzcd}}
% archive sheet number: 67
\scanpage{69}
Donc $T$ est connue (\unsure{à|pr} équiv. près) quand on connaît $S$ et les $r_i:S\longrightarrow(\mathrm{Ens})$ [et on connaît alors \unsure{même|$=$} $T$ sur $\tau^{I}$ \ldots], ou \deleted{de} façon plus intrinsèque, si on connaît \deleted{\unsure{$R\subset$}} \inserted{l'image de $R$ dans} $\underline{\mathrm{Hom}}(S,\mathrm{Ens})=\widehat{S^\circ}$.

\textcircled{3°} On a
\[
S\overset{\sim}{\simeq}\underline{\mathrm{Hom}}_{\Delta}(R,(\mathrm{Ens}))\subset\widehat{R^\circ},
\]
et cette sous-catégorie \uline{contient} $R^\circ$ [car les foncteurs \inserted{cov.} représentables \deleted{\unsure{contra}} de $R$ dans $\mathrm{Ens}$ commutent aux $\varprojlim$ !) donc on a
\[
\boxed{R^\circ\overset{\unsure{j|\partial|\partial\,\xi}}{\subset}S}
\]
et le foncteur \inserted{$\widetilde r:S\longrightarrow(\mathrm{Ens})$} associé à un $r\in\mathrm{Ob}\,R$ est le foncteur qu'il représente \uline{dans} $S$. Ainsi $R$ est connue quand on connaît son image \unsure{ess.|ens.} dans $\underline{\mathrm{Hom}}(S,\mathrm{Ens}]=\widehat{S^\circ}$, car $R\longrightarrow\deleted{\underline{\mathrm{Hom}}}\widehat{S^\circ}$ est le composé $R\xrightarrow{j^\circ}S^\circ\longrightarrow\widehat{S^\circ}$ de deux fonctrs pl. fid.

Ainsi $S=T_{\deleted{\unsure{R}}}(\mathrm{Ens})$ est muni d'une sous-catégorie str. pleine $\Sigma$, ($\approx R^\circ$) [formée des $\varphi\in\mathrm{Ob}\,S$ tels que \deleted{pour} $\forall C\in\mathrm{Ob}\,\Delta$ et tout $\xi\in T(C)$, définissant $\widetilde\xi:C^\circ\longrightarrow S$, \deleted{les} le composé $\inserted{x\mapsto}\mathrm{Hom}(\varphi,\overline{\xi}(x))$ de $C^\circ$ dans $\mathrm{Ens}$ soit représentable]. \inserted{La sous-catégorie} $\Sigma$ est stable dans $S$ par \margin{d'où un foncteur $T(C)\xrightarrow{\unsure{\widetilde\varphi_C|\varphi_C}}C$} $\varinjlim$ (\uline{inductives}) de type $d'$, \deleted{\unsure{de|et} pour tout} On trouve ainsi une équivalence \inserted{explicite} entre $\Sigma^\circ$ et $R$.

\deleted{\textcircled{4°} Soit \deleted{\unsure{$d$}} $\Theta'$ l'\unsure{ens|un|une} de tous les petits diagrammes, $\Delta_0$ \deleted{\unsure{l'une}} la \inserted{2-}catégorie \unsure{de|des} toutes les catégories où les $\varprojlim$ de type $\Theta'$ existent, Supposons que dans la définition de $\Delta$ ne figurent pas des conditions de finitude, de sorte que $\deleted{\unsure{\Delta}}\Delta_0\subset\Delta$. Donc la théorie $T$ \inserted{$=T_R$} 2-repr pour $\Delta$ définit une théorie pour $\Delta_0$, dont nous allons montrer la 2-représentabilité. \deleted{Soit $C\in\mathrm{Ob}\,\Delta_0$, on a} \uline{Si} c'est représentable, on}
\[
\deleted{\underline{\mathrm{Hom}}_{\Delta_{\unsure{0|4}}}(R,C)\simeq\underline{\mathrm{Hom}}_{\varinjlim}(R^\circ,C^\circ)^\circ\unsure{\simeq|\sim}}
\]
\deleted{sait que la catégorie $R_0$ qui représente \deleted{\unsure{sera|dans}} une catégorie avec \deleted{$\varinjlim$} $\varprojlim$ quelconques, \deleted{\unsure{muni de}} etc}
% archive sheet number: 68
\scanpage{70}
% archive sheet number: 69
\uline{Exemples}

\textcircled{1} $d'=d''=\varnothing$. \deleted{Alors} a) Prenant $\Delta=(\mathrm{Cat}_d)$, les Pbs 1) et 2) ont une solution évidente resp. bien connue. Les théories Algébriques sont exactement celles qui sont représentées par des \uline{petites} catégories. Les $I\longrightarrow R$ générateurs sont uniquement les ess. surjectifs.

b) \uline{Condition} d'exactitude. On peut par exemple prendre \uline{$\Delta=(\text{groupoïdes})$}.
\margin{\diagram{$\underline{\unsure{\rho|p}}=\left[\,\begin{tikzcd}[ampersand replacement=\&,column sep=small]
\overset{a}{\bullet}\arrow[r,"u"] \& \overset{b}{\bullet}
\end{tikzcd}\,\right]$}
\unsure{$u$ rendu un|on rend $u$ un} isom}
Les \deleted{\unsure{cas}} Pb 1 a une solution connue (groupoïde fondamental d'une catégorie), le Pb 2) est trivial. Les théories Alg. sont celles repr. par des petites \deleted{catégories} \inserted{groupoïdes}. Les $I\longrightarrow R$ générateurs sont les ess. surj.

c) Prenons $\Delta=(\text{catégories ordonnées})$. Pb 1 a une solution évidente (\deleted{\unsure{catégori|catégorie}} ens. préordonné associé à une catégorie), Pb 2 une solution facile.
\margin{\diagram{$\underline{\unsure{\rho|p}}=\left[\,\begin{tikzcd}[ampersand replacement=\&,column sep=small]
\underset{a}{\bullet}\arrow[r,shift left,"u"]\arrow[r,shift right,"v"'] \& \underset{b}{\bullet}
\end{tikzcd}\,\right]$}
\deleted{\unsure{pour|que}} \unsure{ici|là} on rend $u$ égal à $v$}
Les théories algébriques sont \inserted{celles} repr. par petits ens. \unsure{(pré)ordonnés|(pré), ordonnés}. Les $I\longrightarrow R$ générateurs sont ess. surjectifs.
\scanpage{71}
\textcircled{2} $d'=$ tous les ``petits'' diagrammes, \inserted{$d''=\varnothing$.} Pas de conditions d'exactitude.

\uline{Pb 1} $I$ une petite catégorie, on a pour $C\in\mathrm{Ob}\,\Delta$
\[
\underline{\mathrm{Hom}}(I,C)\simeq\underline{\mathrm{Hom}}(I^\circ,
\underset{\substack{\text{\unsure{avec|sans}}\\\varinjlim\ \text{quelc.}}}{C^\circ})^{\circ}
\simeq\underline{\mathrm{Hom}}_{\Delta^\circ}(\widehat{I^\circ},C^\circ)^\circ
\simeq\underline{\mathrm{Hom}}_{\Delta}(\widehat{I^\circ}{}^\circ,C)
\]
donc $I$ a la $\Delta$-enveloppe $\widetilde I=\underline{\mathrm{Hom}}(I,\unsure{C|e})^\circ$ (qui est l'opposé d'un topos). Notons que $I\longrightarrow\widetilde I$ définit les objets de $I$ comme famille \uline{cogénératrice} (donc $I^\circ\longrightarrow\widetilde I^\circ$ comme famille génératrice).

\uline{Pb 2} $C\in\mathrm{Ob}\,\Delta$, $\Sigma\subset\mathrm{Fl}\,C$ [$\Sigma$ petit]. Rendre inversibles les éléments de $\Sigma$ ? \deleted{Supposons que} Or $C\Sigma^{-1}$ (au sens de Cat) \inserted{est défini, montrer que c'est une solution,}. On \uline{peut supposer} que $\Sigma$ est stable par composition, par chgt de base, contient les isom, et \inserted{tel} que $vu\in\Sigma$, $v\in\Sigma\Longrightarrow u\in\Sigma$. Mais alors $\Sigma$ donne lieu à un ``calcul de fractions \unsure{à g|: g}'', et il est connu que \inserted{dans $C\Sigma^{-1}$} les \unsure{dp|types} pour les petites $\varprojlim$ qui existent dans $\unsure{\mathscr C|\mathbf C|C}$ existent dans $C\Sigma^{-1}$, et $C\longrightarrow C\Sigma^{-1}$ y commute.
% slanted
\margin{$\Sigma$ ne sera plus petit après saturation !}
Il reste à examiner les produits \inserted{(infinis)}. Pour ceci, on peut supposer aussi que $\Sigma$ est stable par produits \uline{infinis} de morphismes. Mais alors on voit tout de suite que les produits de $C$ sont des produits de $C\Sigma^{-1}$. OK !
% slanted
\margin{Mais il faut \deleted{Si $C$ est} vérifier que la catégorie obtenue est une $\mathcal U$-catégorie.}

Notons que si $I$ est une famille ``cogénératrice \unsure{pour mon strict|pour mon stricts|par mono strict}'' dans \deleted{Pour tout} $C$ [pour tt $X\in\mathrm{Ob}\,C$, $X\xrightarrow{\sim}\varprojlim_{X^{\prime}/C}X'$], alors il en est encore de \unsure{même|$=$} dans $C\Sigma^{-1}$. \deleted{Donc on voit}

Supposons qu'on ait un \inserted{\unsure{premier|pseudo-}} syst. inductif filtrant $(R_\alpha)$ dans $\Delta$, avec $I\longrightarrow R_\alpha$ cogénérateur. Soit $R=\underset{\alpha}{\varinjlim}_{\Delta}R_\alpha$. Alors $I\longrightarrow R$ est-il cogénérateur ? Soit $R'=\underset{\alpha}{\varinjlim}_{\mathrm{Cat}}R_\alpha$. \deleted{Dans $R'$, les $\varprojlim$ \unsure{prennent|proviennent} de syst.} Donc $R$ s'obtient à partir de $R'$ en rendant inversibles certaines flèches, donc $\mathrm{Ob}\,R^{\unsure{\prime\prime|\prime}}=\bigcup_{\alpha}\mathrm{Im}\,R_\alpha$. De ceci on déduit aussitôt ce qu'on a dit.
% slanted
\margin{\uline{Mais attention} le fait d'être strictement cogénérateur \uline{n'est} \uline{pas} transitif}

Il faut aussi étudier le cas où on ``rajoute des flèches'' à $R$ muni de $I$, d'où $R\longrightarrow R'$. Ce \deleted{\unsure{dernier}} \inserted{$\Delta$-foncteur} est ess. surjectif, et on gagne encore !
% archive sheet number: 70
\scanpage{72}
\uline{Th.} Les théories algébriques pour $\Delta$ sont celles qui sont 2-représentées par des $R\in\mathrm{Ob}\,\Delta$ admettant des \uline{petites familles \inserted{str.} cogénératrices}. Pour que $I\xrightarrow{\rho_0}R$ \deleted{\unsure{avec petite image dans|ayant petite image dans}} soit \inserted{$\Delta$-}basique, il f. et s. que son image soit \inserted{str. $\Delta$-}cogénératrice (c'est plus faible que \inserted{str.} cogénératrice !).
% slanted
\margin{passe bien dans la construction transfinie \unsure{si on prouve|d'en prouver} stabilité par $R[H]_\Delta$}

La nécessité dans la première assertion est claire par ce qui précède. Pour la suffisance, \deleted{on peut supp.} elle résulte de la deuxième assertion. Or soit $I'$ la catégorie \inserted{mun. d'un foncteur \uline{pl. fidèle} $I'\xrightarrow{\rho}R$,} telle que $\mathrm{Ob}\,I'=I$, $\mathrm{Ob}\,\rho=\varphi_0$. \unsure{Considérons|Considérons. On|\deleted{Considérons} On} \deleted{alors} \unsure{Sait|sait} que $T_{\widetilde I'}$ est une théorie algébrique de base $I$. Or $\widehat{I'^\circ}\longrightarrow R^\circ$ prolongeant $\rho:I'\longrightarrow R$ est ess. surjectif \inserted{et surjectif sur les Hom,} donc $\widetilde I'\longrightarrow R$ aussi. Donc il \deleted{suffit de} reste à rendre \deleted{\unsure{égales}} des \deleted{\unsure{iso}} flèches, donc rendre certaines flèches des iso.

\uline{Remarques} \textcircled{1} On \deleted{précise} \inserted{obtient} ici \deleted{à} la construction de tte théorie algébrique en les étapes suivantes : 1°) \unsure{choix|élts|objets} de base $I$ 2°) rajouter des flèches 3°) axiomes d'égalité sur des \inserted{composés de} flèches (on a maintenant $T_{\widetilde I'}$), 4°) \deleted{axiomes d'égalité sur d'autres \inserted{composés de} flèches, ou encore} \inserted{ou encore rendre inversibles des flèches}

On peut d'ailleurs condenser 3° et 4°) en une seule étape, \deleted{\unsure{3}} (imposer des axiomes successivement \unsure{revient à|rev. à} le faire tjrs en une seule fois, grâce au calcul de fractions) donc on donne $R$ \deleted{par} comme catégor. de fractions de $\widetilde I'=\widehat{I'^\circ}{}^\circ$ (d'un opposé de topos) par un ens. de flèches satisfaisant les conditions de saturation déjà dites.

\textcircled{2} Soit $R\in\mathrm{Ob}\,\Delta$. Les catégories quotients $R'\in\mathrm{Ob}\,\Delta$ déduites de $R$ \deleted{par} rendant inversibles des flèches \deleted{sont exactement celles qui} \inserted{correspondent 1--1 aux} ens. $\Sigma$ de flèches \deleted{$\Sigma$} telles que

a) $\Sigma$ stable par composition, par changement de base, contient les iso

b) $\Sigma$ stable par produits quelconques.

\textcircled{3} Supposons que $R$ ait une petite famille cogénératrice $I$. Alors \deleted{les} $\underline{\mathrm{Hom}}_{\Delta}(R,\mathrm{Ens})\simeq R^\circ$ [les foncteurs de $R$
% archive sheet number: 71
\scanpage{73}
dans $(\mathrm{Ens})$ qui commutent aux $\varprojlim$ sont les foncteurs représentables) i.e. si $T=T_R$, on a $R=S^\circ$, où $S=T(\mathrm{Ens})$. \deleted{Cela précise} \deleted{Donc, \unsure{pour les}} Ou encore $R^\circ=S$. Donc :

\uline{Cor} \deleted{les} \inserted{2-cat des} structures algébriques définissables par $\varprojlim$ (sans axiome d'exactitude) \deleted{correspondent aux catégories} est 2-équiv. par $T\longmapsto T(\mathrm{Ens})$ à celle des $\mathcal U$-catégories $S$ avec $\varinjlim$ quelc. et \inserted{petites} \deleted{familles} \inserted{sous-cat} génératrices par \unsure{épis.|épi.} stricts [i.e. telles que $S\longrightarrow\widehat{I^{\unsure{a|\circ}}}$ soit conservatif]. La donnée d'une base \inserted{de $T$} équivaut à celle d'une famille d'objets cogénératrice de $S$. Pour toute $C\in\mathrm{Ob}\,\Delta$, on a
\[
T(C)\simeq\underline{\mathrm{Hom}}_{\Delta}(S^\circ,\unsure{\mathscr C|\mathbf C|C})\deleted{\unsure{\mathrm{Ens}}}
\]

L'équivalence est obtenue en notant que le deuxième mb est formé des foncteurs tels que $\forall x\in C$, le composé avec $\underline{\mathrm{Hom}}_{C}(x,-)$ est représentable, donc le deuxième mb \deleted{est} peut s'interpréter comme les foncteurs $S^\circ\times C^\circ\longrightarrow(\mathrm{Ens})$ qui sont représentables pour chaque argument fixé, ou encore les foncteurs $C^\circ\longrightarrow S$ qui sont tels que pour tout $Y\in S$, son composé avec $\mathrm{Hom}(Y,-)$ soit représentable.

Or on a
\[
T(C)\longrightarrow\underline{\mathrm{Hom}}(C^\circ,\underbrace{T(\mathrm{Ens})}_{\substack{\Vert\\S}})=\underline{\mathrm{Hom}}(\unsure{C^\circ|\mathscr C^\circ},S)
\]
déduit de $C^\circ\longrightarrow\underline{\mathrm{Hom}}_{\Delta}(C,\mathrm{Ens})\subset\widehat{C^\circ}$ induisant $C^\circ\longrightarrow\underline{\mathrm{Hom}}(T(C),T(\mathrm{Ens}))$ i.e. $T(C)\longrightarrow\deleted{\unsure{\mathrm{Hom}|T(C)}}\underline{\mathrm{Hom}}(C^\circ\ T(\mathrm{Ens}))$.

\inserted{\deleted{\unsure{Ens}} \textcircled{4} Comme $S=T(\mathrm{Ens})$ a des $\varprojlim$ quelc, et $S\cong R^\circ$, $R$ a aussi des $\varinjlim$ quelc, donc $S$ et $R$ ont toutes les $\varinjlim$ et $\varprojlim$}

\textcircled{\deleted{4}5} Soit un foncteur $\varphi:R\longrightarrow R'$ entre objets de $\Delta$. Si $R$ a une petite famille cogénératrice, alors $\varphi$ est un $\Delta$-foncteur, i.e. commute aux $\varprojlim$, sss il admet un adjoint à gauche $\psi$
\[
\mathrm{Hom}(y',\varphi(y))\simeq\mathrm{Hom}(\psi(y'),x)
\]
% slanted
\margin{ainsi les \inserted{$\Delta$-}foncteurs $\varphi:R\longrightarrow R'$ (i.e. $T'\longrightarrow T$) correspondent aux foncteurs $\psi:R'\longrightarrow R$ (i.e. $\psi^\circ:S'\longrightarrow S$) qui ont un adjoint à droite (\unsure{Corresp à|corresp. à} \unsure{$\varphi^\circ$})}
Alors $\varphi$ est un passage à une sous-cat. \unsure{mise|univers.} en fractions sss \deleted{$R'$ est une sous-catégorie} $\psi$ est pl fidèle, \inserted{(et alors $R'$ a une p. famille cog.)} Ainsi les \uline{sous-}\deleted{catégories}\inserted{th$^{\text{é}}$}théories \inserted{Alg.} de la th. algébrique $R$ correspondent \deleted{aux sous-} à \uline{certaines} sous-catégories pleines de $R$, savoir celles qui ont des $\varprojlim$ et telles que les foncteurs d'inclusion $\psi$ ait un adjoint à \deleted{\unsure{g}} droite. En termes de $S=R^\circ=T(\mathrm{Ens})$, et \inserted{$S'=R'^\circ=T'(\mathrm{Ens})$} \deleted{ce sont les sous-catégo-}
% archive sheet number: 72
\scanpage{74}
\deleted{ries} $\psi^\circ:\overset{\substack{\inserted{S'=R'^\circ}\\\Vert}}{T'(\mathrm{Ens})}\longrightarrow\overset{\substack{\inserted{S=R^\circ}\\\Vert}}{T(\mathrm{Ens})}$ est simplement le foncteur associé à $T'\longrightarrow T$ déduit de $\varphi$. Ce foncteur doit donc commuter aux $\varprojlim$ quelc. [\unsure{de ce|ce|donc ce} qui exprime simplement que $\varphi$ commute aux $\varinjlim$, \inserted{\deleted{\unsure{et on|alors}} or c'est un adjoint à g.}] Donc les sous-théories de $T$ correspondent aux sous-catégories pleines \inserted{$S'$} de $S$, \deleted{\unsure{admettant un|et permettant de|admettant des}} \inserted{[stables par} $\varprojlim$, ce sera automatique par les autres conditions] admettant des $\varinjlim$ quelconques (mais pas néc. stables par ces dern) et telles que le foncteur d'inclusion $S'\overset{i\,(=\psi^\circ)}{\hookrightarrow}\unsure{S|S^\circ}$ ait un adjoint à \uline{gauche} \deleted{$\unsure{\varphi}$} $j\ (= \varphi^\circ)$
\[
\mathrm{Hom}_{S}(x,i(y'))\simeq\mathrm{Hom}(jx,y')
\]

\deleted{Alors on a des foncteurs, \illegible{1} adjts de \unsure{oubli que|oubliques}} \deleted{[Mais en général on (Structure} ($S'$ (foncteur $S'$-structure ``libre'' définie par une $S$-structure) Mais en général, il n'y a pas moyen de définir pour tte $C\in\mathrm{Ob}\,\Delta$ un foncteur $T(C)\longrightarrow T'(C)$ qui \deleted{\unsure{en|les}} redonne $j$ \deleted{\unsure{comme $C$}} pour $C=(\mathrm{Ens})$ [\uline{NB} Cependant, si $C$ est un \uline{topos}, et si $T,T'$ et $\varphi$ se définissent au stade des $\varprojlim$ finies, c'est Ok \ldots]

\uline{NB} La 2-catégorie des théories algébriques pour $\Delta$ est 2-équivalente à la 2-catégorie des cat. $S$ qui sont stables par \inserted{petites} $\varinjlim$ \deleted{quelc.} et qui ont une petite sous-catégorie str. génératrice, avec les foncteurs qui commutent aux $\varinjlim$ (mais en \unsure{renversant|inversant} les flèches pour les morphismes ?). \deleted{Ce sont aussi les catégories telles que les \inserted{contrav.} foncteurs qui transforment $\varinjlim$ quelc en limites \unsure{proj.}}
% archive sheet number: 73
\scanpage{75}
\textcircled{3} \deleted{$d'=$ \unsure{tous les|toutes les}} $d'$ \deleted{formé de diagrammes finis} \inserted{quelconque}, $d''=\varnothing$, \uline{pas de conditions d'exactitude}.

\deleted{\unsure{Pb 1} On suppose que} Soit $d'_0$ l'ens de tous les diagrammes finis, $d'_1$ l'ens de tous les petits diagrammes, et $\Delta\supset\Delta_0\supset\Delta_1$ les inclusions correspondantes.

\uline{Pb 1} $\Delta$-Enveloppe \inserted{$\rho_{\Delta}(I)$} d'une catégorie $I$. Soit $\overline I$ \deleted{$\unsure{\rho(I)}$ \illegible{1}} $\simeq\widehat{I^\circ}{}^\circ$ la $\Delta_1$ enveloppe, et soit $\widetilde I$ la \deleted{sous}-catégorie pl. de $\overline I$ formée des $\varprojlim$ de diagr. $d'$ d'objets de $I$ [NB $I\longrightarrow\overline I$ pl fid !].

Je dis que cela marche, \uline{pourvu} qu'on sache que $\widetilde I$ est stable par $d'$ (c'est le cas si $d'=$ ens. des diagrammes finis discrets, ou $d'=$ ens. de tous les diagrammes finis \inserted{(!)} \ldots) sinon on doit itérer et prendre \deleted{la lim} pour $\widetilde I$ la plus petite sous catégorie \inserted{str} pleine de $\overline I$ contenant $I$ et stable par $\varprojlim$ de type $d'$. En effet, on a pour $C\in\mathrm{Ob}\,\Delta$
% slanted
\margin{Pb 1 et Pb $1'$ Cela marche pour \uline{tout} ens. $d'$ de petits diagrammes (pas néc. finis)}
\[
\deleted{\underline{\mathrm{Hom}}_{\deleted{\unsure{\Delta|\mathrm{Cat}}}\mathrm{Cat}}(I,\unsure{\mathscr C|C})\longrightarrow\underline{\mathrm{Hom}}_{\Delta_1}(\overline I,\overline C),}
\]
\deleted{\deleted{et} \inserted{foncteur (pl. fidèle dont)} l'image ess. est formée de(s) foncteurs $\overline I\xrightarrow{\overline\varphi}\overline C$ tels que $\overline\varphi(I)\subset\deleted{\unsure{\mathrm{Im}}} C$ (ess.) Mais \deleted{alors} $\overline\varphi^{-1}(C)$}

Soit $\rho(C)$ la $\Delta$-$\Delta_1$ enveloppe de $C$ i.e. la $\Delta_1$-catégorie avec $\Delta$-foncteur $C\xrightarrow{i}\rho(C)$ \uline{universel} [c'est une catégorie \deleted{\unsure{quot.|pleine}} de fraction de $\overline C$]. \deleted{\unsure{Alors supposons que}} \inserted{Admettons que} \deleted{$I\longrightarrow C$ se prolonge} ce foncteur est pleinement fidèle (Pb auxiliaire) \deleted{En effet, \unsure{c'est} les foncteurs} Alors les foncteurs $\varphi:I\longrightarrow C$ \deleted{définist, \unsure{via i$\varphi$|à $\varphi$}} \unsure{``s.a''|``s.a. à''} les foncteurs $\psi=i\varphi:\widetilde I\longrightarrow\rho(C)$ qui envoient $I$ dans l'image essentielle de $C$ par $i$, ou encore les foncteurs $\overline\psi:\overline I\longrightarrow\rho(C)$. tels que $\overline\psi(I)$ soit contenue dans l'image essentielle de $\unsure{I|i|R|C}$ ; de tels $\overline\psi$ envoient $\widetilde I$ dans $C$ d'où $\underline{\mathrm{Hom}}_{\mathrm{Cat}}(I,C)\longrightarrow\underline{\mathrm{Hom}}_{\Delta}(\widetilde I,C)$, et on a un foncteur ``restriction'' en sens inverse. Il est clair que les composés dans les deux sens sont l'identité, OK.

\uline{Pb $1'$} \uline{Ou} Pb \uline{auxiliaire} Si $C\in\mathrm{Ob}\,\Delta$, montrer que $C\longrightarrow\rho(C)$ \uline{est} \uline{pl.} \uline{fidèle}. Je dis qu'on a
\[
\rho(C)=\underline{\mathrm{Hom}}_\Delta(C,\mathrm{Ens})^\circ=T_C(\mathrm{Ens})^\circ,
\]
% archive sheet number: 74
\scanpage{76}
le foncteur $C\longrightarrow\rho(C)$ \inserted{$\subset\underline{\mathrm{Hom}}(C,\mathrm{Ens})^\circ=\widehat{C^\circ}^{\,\circ}$} donné \deleted{par} en associant à tt objet de $C$ le foncteur covariant qu'il représente (alors il est clair que $C\longrightarrow\rho(C)$ est pl. fidèle !). On a en effet, \deleted{si $D$ est une catégorie où les $\varprojlim$ quelc. sont \unsure{repr.|requ.}, et si $\rho(C)$ est défini comme la $\Delta$-$\Delta_1$ enveloppe $\underline{\mathrm{Hom}}_{\unsure{\Delta_1|\Delta}}(C,D)\simeq\underline{\mathrm{Hom}}_{\Delta_1}(\rho(C),D)\simeq\unsure{\mathrm{H}}$}
\[
\begin{aligned}
\rho(C)&\cong T_{\rho(C),\Delta_1}(\mathrm{Ens})^\circ\\
&\cong T_{C,\Delta}(\mathrm{Ens})^\circ=\underline{\mathrm{Hom}}_\Delta(C,\mathrm{Ens})^\circ\mathrel{\therefore}\mathrm{OK}.
\end{aligned}
\]

\uline{Pb 2} $C\in\Delta$, $\Sigma\subset\mathrm{Fl}$, construire $C\Sigma^{-1}$ \uline{dans $\Delta$}, et de m\textsuperscript{me} si on se donne un ens. de couples de flèches qu'on veut égaliser. \deleted{Tout marche bien si $d'$} Soit $\rho(C,\Sigma)$ la \inserted{$\Delta_1$-}catégorie de \deleted{quotient}\inserted{fractions} de $\rho(C)$, définie par $\Sigma$, et soit $C'$ la plus petite sous-\inserted{$\Delta$-}catégorie \inserted{str. pleine} de $\rho(C,\Sigma)$ \deleted{qui} telle que $C\longrightarrow\rho(C,\Sigma)$ se factorise par $C'$, \inserted{i.e. telle que $C'$ contienne la sous-catégorie \uline{pleine} $C\Sigma^{-1}$.}

% slanted
\margin{Si $\Sigma$ contient les diagrammes des produits fibrés et des noyaux, alors il semble que $C\Sigma^{-1}$ \deleted{\unsure{pourra}}\inserted{convienne}, voir si pour tt $\delta\in d'$, les produits fibrés de type $\mathrm{Ob}\,\delta$ existent dans $C'$, i.e. \deleted{\unsure{le}} $d'$ contient le diagramme \diagram{$\begin{tikzcd}[ampersand replacement=\&,column sep=tiny,row sep=small] {} \arrow[dr] \& {} \arrow[d] \& {} \arrow[dl]\\ \& \bullet \& \end{tikzcd}$} \unsure{où l'un au moins|ou l'un analogue} \unsure{sur un ens.|est} $\times$ $\mathrm{Ob}\,\delta$. OK. Si p. ex. $d'$ est l'ens. de tous les diagrammes \uline{finis}, mais non si c'est l'ens. des diagrammes finis \uline{discrets}.}

Soit $D\in\mathrm{Ob}\,\Delta$, alors les $\Delta$-foncteurs $\varphi:C\longrightarrow D$ correspondent aux $\Delta_1$-foncteurs $\psi:\rho(C)\longrightarrow\rho(D)$ tels que $\psi$ envoie $C$ dans (la ss-cat. pleine) $D$, et les $\Delta$-$\Sigma$-foncteurs ``sont'' \deleted{tels que} les foncteurs $\psi':\rho(C,\Sigma)\longrightarrow\rho(D)$ \deleted{tels que} qui envoient $C$ dans $D$, \deleted{ou encore} ce qui revient au m\textsuperscript{me}, $C'$ dans $D$. Donc
\diagram{$\begin{tikzcd}[ampersand replacement=\&]
\underline{\mathrm{Hom}}_{\Delta,\Sigma}(C,D) \arrow[r,"\alpha"] \& \underline{\mathrm{Hom}}_\Delta(C',D) \arrow[l,bend left,"\beta"]
\end{tikzcd}$}
, et on a le foncteur restriction en sens inverse. Il est clair que les deux composés sont \deleted{égaux}\inserted{$\simeq$} les foncteurs identiques, trivial pour $\beta\alpha$, pour $\alpha\beta$ on note que \deleted{\unsure{tout d'abord} $\beta$ est pl. fidèle,} $C'$ est déduit de $C\Sigma^{-1}\underset{\mathrm{pl.}}{\subset}\rho(C,\Sigma)$ en rajoutant transfiniment des $\varprojlim$ de type $d'$, et par suite que $\underline{\mathrm{Hom}}_\Delta(C',D)\longrightarrow\underline{\mathrm{Hom}}_{\unsure{\Delta|\subset\Delta|\mathrm{Cat}}}(C,D)$ est pl. fidèle (car $\underline{\mathrm{Hom}}_\Delta(C',D)\longrightarrow\underline{\mathrm{Hom}}_{\mathrm{Cat}}(C',D)$ l'est \ldots).

\uline{NB} Dans les deux cas envisagés (Pb 1 et Pb 2) la catégorie résultante (de $I$, resp. de $C$) se déduit de la sous catégorie pleine de départ en rajoutant des $\varprojlim$ de type $d'$ (transfiniment au besoin). On trouve sauf erreur :

% slanted
\margin{il faudrait encore regarder le cas des $\unsure{\lim_\alpha|\varinjlim_\alpha|\varprojlim_\alpha}$ et du \unsure{rajoutage|v. multiple} de flèches}

\uline{Th.} Les théories algébriques de type $\Delta$ sont celles qui sont 2-représentées par des $R\in\mathrm{Ob}\,\Delta$ ayant une petite sous-catégorie \inserted{$I'$} \deleted{qui} $d'$-\inserted{co}génératrice [i.e. \deleted{tt objet de} $R$ est la réunion d'une suite transfinie de sous-catégories str. pleines de $R$ déduites en rajoutant des $\varprojlim$ de type $d'$]. Un foncteur $I\longrightarrow R$ \inserted{($I$ discret)} est basique sss il est $d'$-cogénérateur.
% archive sheet number: 75
\scanpage{77}
On est encore ramené à prouver ce dernier énoncé. On procède comme dans 2°, en introduisant la sous-cat. pleine $I'$ de $R$ telle que $\mathrm{Ob}\,I'=I$, d'où $T_R\longrightarrow T_{\widetilde I'}$ \uline{pl. fidèle}. Regardons $\widetilde I'\longrightarrow R$, il est \inserted{ess.}\deleted{surjectif sur les objets et} surjectif sur les Hom, \deleted{\unsure{donc}} est-il ess. surjectif ? OK. Si \deleted{le foncteur} $d'$ \deleted{\unsure{contient}} satisfait aux conditions écrites en marge plus haut, car alors on trouve un \inserted{$\Delta$-}foncteur pl. fidèle $\underbrace{\widetilde I'\Sigma^{-1}}_{\substack{\text{catégorie}\\\text{de fractions}\\\text{habituelle}}}\longrightarrow R$, \deleted{\unsure{dont}} qui doit donc être ess. surjectif. Donc on trouve

\diagram{$\begin{tikzcd}[ampersand replacement=\&,column sep=small,row sep=small]
\& \widetilde I' \arrow[ddl] \arrow[dl,dash] \arrow[dr,dash] \& \deleted{\unsure{I}}\\
I' \arrow[rr] \arrow[d,hook,"\text{pl. fid}"'] \& \& I' \arrow[d,"\substack{\text{ess. surj.}\\\text{et surjectif sur Hom}}"]\\
R \& \unsure{\subset|\longleftarrow|\hookleftarrow} \& \rho(R)
\end{tikzcd}$}

$I'$ coengendre strict\textsuperscript{mt} $R$ qui coengendre str. $\rho(R)$ donc $I'$ coengendre str. $\rho(R)$ (?)

\uline{Corollaire} \inserted{Si $d'$ satisfait aux conditions envisagées (p. ex. $d'=$ ens. des diagrammes finis)} On peut construire les théories algébriques pour $\Delta$ en trois pas, comme pour $\Delta_1$ : on prend une catégorie $I'$ ayant $I$ comme ens. d'objets (par 1° et 2°) et prend une catégorie de fractions de $\widetilde I'$.

% Two vertical strokes beside this paragraph.
Mais la m\textsuperscript{me} conclusion est valable si p. ex. $d'$ est formé de diagrammes discrets, stable par passage à un sous-diagramme \unsure{(|\{} \deleted{\unsure{diagrammes}}

Néanmoins, dans le cas d'un $d'$ général (m\textsuperscript{me} formé de diagrammes finis) il semble qu'il faille une construction récurrente transfinie

\deleted{\unsure{Th|Pb} \unsure{Question} \inserted{\uline{Question}} Supposons $d'$ petit. Alors une $R$ \deleted{telle que $R$} admettant une petite ss-catégorie $d'$-cogénératrice (donc \deleted{2-rep} correspondant à une théorie algébrique pour $\Delta$) est petite, car on \unsure{a}\ldots}
% archive sheet number: 76
\scanpage{78}
4°) Cas $d'=$ ens des diagrammes finis, $d''=\varnothing$, pas de conditions d'exactitude.

Comme $d'$ est petit, les $R\in\mathrm{Ob}\,\Delta_0$ qui 2-représentent des th. algébriques sont les $R$ ess. petites (stables par $\varprojlim$ finies). \deleted{\unsure{Donnons-nous|Ramenons} \unsure{donc} $R$, lui \unsure{associons} $\unsure{\rho(R)}=\unsure{R^\circ}$ \unsure{Hom}\, $\unsure{(R,\mathrm{Ens})}$}

\[
\left\{\begin{aligned}
&T\quad\text{théorie algébrique pour }\Delta\\
&R\in\mathrm{Ob}\,\Delta\quad\text{qui représente }T\\
&S=T(\mathrm{Ens})=\underline{\mathrm{Hom}}_{\mathrm{lex}}(R,\mathrm{Ens}))\\
&\Sigma\subset S\quad\text{(structures représentables)}
\end{aligned}\right.
\]

Tout se détermine mutuellem\textsuperscript{t}

\[
\begin{aligned}
R:\qquad&T=T_R:C\longrightarrow\underline{\mathrm{Hom}}_{\mathrm{lex}}(R,C)\\
&S=\deleted{\unsure{T(\mathrm{E})}}\;\underline{\mathrm{Hom}}_{\mathrm{lex}}(R,\mathrm{Ens})\\
&\Sigma\cong R^\circ\subset S\quad\text{par foncteurs représ.}\\[2pt]
\Sigma:\qquad&\unsure{R|\ell}=\Sigma^\circ\\
&T:\unsure{R|C}\longmapsto\underline{\mathrm{Hom}}_{\mathrm{lex}}(\Sigma^\circ,C)\\
&S=\underline{\mathrm{Hom}}_{\deleted{\unsure{\mathrm{en}}}\mathrm{lex}}(S^\circ,C),\quad\text{contenant }\Sigma\text{ comme sous-catégorie pleine.}\\
&\deleted{\unsure{\text{ss-cat. repr.}}}
\end{aligned}
\]

\deleted{\unsure{Reste à voir une}} Autre façon de décrire $S$ de $R$ ou $\Sigma$

% slanted
\margin{cf SGA 4 I 8.7.5 où (i) $\unsure{\Sigma\in S|\Sigma_i\in S}$ où $\unsure{F|\Sigma}\in\underline{\mathrm{Hom}}(\unsure{\Sigma^\circ},\mathrm{Ens})$ (ii) $\Sigma$ strictement génératrice (iii) $\Sigma$ stable par petites $\varinjlim$, $\Sigma$ petite, le foncteur \unsure{:} ex. à dr.}

\uline{Lemme} \deleted{\unsure{On}} L'inclusion $\Sigma\longrightarrow S$ définit une \uline{équivalence}
\[
\mathrm{Ind}(\Sigma)\longrightarrow S
\]

\margin{SGA 4 I 8.9.5.}
[ En effet, \deleted{\unsure{est}} on sait que $\Sigma\longrightarrow S$ est universel pour des foncteurs $\Sigma\longrightarrow D$ ex. à dr, où $D$ a des $\varinjlim$ quelcon. Or montrons que $\Sigma\longrightarrow\mathrm{Ind}(\Sigma)$ est \deleted{\unsure{ident}} itou. \inserted{On sait que $\mathrm{Ind}(\Sigma)$ a des $\varinjlim$ finies (donc quelc.) car $\Sigma$ en a, et $\Sigma\longrightarrow\mathrm{Ind}(\Sigma)$ y comm.} \deleted{En effet} les foncteurs \uline{quelc.} $\Sigma\longrightarrow D$ \deleted{\unsure{sont les}}\inserted{correspondent}
\margin{SGA 4 I 8.7.3}
aux foncteurs $\mathrm{Ind}(\Sigma)\longrightarrow D$ qui commutent aux $\varinjlim$. Ceux qui commutent aux $\varinjlim$ finies (donc quelconques) sont ceux qui sur $\Sigma$ y commutent [nécessité triviale, pour la suffisance, on \deleted{note que \unsure{cela} revient aux \unsure{copr.} finies, initiaux, sommes, conoyaux,} \inserted{note que par I 8.9.8 $\mathrm{Ind}(\Sigma)\longrightarrow\mathrm{Ind}(D)$ \deleted{commute} est exact à droite,} or $\mathrm{Ind}\,D\longrightarrow D$ l'est aussi (c'est un adjoint à g, qui commute à \uline{ttes} limites inductives) ok.

\uline{Corollaire} \deleted{Pour tout} \inserted{Soit $X\in S$. Pour que} $X\in\Sigma$, \inserted{il f. et s. que} le foncteur $\mathrm{Hom}(X,-)$ sur $S$ commute aux $\varinjlim$ filtrantes.

% slanted
\margin{\uline{NB} Marche pour tte catégorie karoubienne ! cf aussi SGA 4 I 8.7.8}

C'est évidemment nécessaire, par construction des Hom dans $\mathrm{Ind}(S)$.

C'est suffisant, car \deleted{si $X\in S$,} On aura $X=\varinjlim_i X_i$ ($X_i\in\mathrm{Ob}\,\Sigma$), filtrante, donc le morphisme identique se factorise par un des $X_{\unsure{j|i}}$, donc on a un $X_i\longrightarrow X$ ayant une section, donc $X$ est une \deleted{\unsure{image}} image projecteur, or $\Sigma$ est \deleted{donc} stable \inserted{pour ça} (c'est un noyau) \deleted{et $\Sigma\longrightarrow S$, est exact à dr, donc} OK. $]$
% archive sheet number: 77
\scanpage{79}
Donc $\Sigma$ se réconstitue à partir de $S$ comme l'\deleted{\unsure{ens. des|ens. des ob.}} sous-catégorie (str.) pleine des objets ``de présentation finie'' !

\[
\begin{aligned}
S:\qquad&\Sigma=S_{pF}\subset S\\
&R=\Sigma^\circ\\
&T:C\longmapsto\underline{\mathrm{Hom}}_{\mathrm{lex}}(S_{pF}^{\,\circ},\unsure{C|e})
\end{aligned}
\]

\[
\begin{aligned}
T:\qquad&S=T(\mathrm{Ens})\\
&\Sigma=S_{pF}=T(\mathrm{Ens})_{Fp}\\
&R=\Sigma^\circ
\end{aligned}
\]

\uline{Corollaire} Pour qu'une \deleted{\unsure{th.|théorie}} $\Delta_1$-théorie \inserted{$T$} (pour $\varprojlim$ quelc) provienne d'une théorie définissable par $\varprojlim$ finie, il f. et s. que $S=S_T=T(\mathrm{Ens})$ soit \deleted{\unsure{engendré}} strictement engendrée par \deleted{ses} \inserted{sa sous-catégorie pleine $\Sigma$ \unsure{de ses|des}} objets de \unsure{prés.|Pr.} finie, ou encore que $S\longrightarrow\deleted{\unsure{\mathrm{Ens}}}\widehat{\Sigma}$ soit pleinement fidèle. (Alors $\Sigma^\circ$ est la catégorie classifiante \deleted{pour} la théorie sur $\Delta_0$ qui donne naissance à $T$).

\uline{Corollaire} Le 2-foncteur de restriction à $\Delta_1$ des théories algébriques sur $\Delta_0$ est 1-fidèle, i.e. pour \deleted{\unsure{la}} deux théories $T,T'$ sur $\Delta_0$, on a
\[
\underline{\mathrm{Hom}}_{\mathrm{Th}_{\deleted{\unsure{\Delta_0}}\Delta_0}}(T,T')\longrightarrow\underline{\mathrm{Hom}}_{\mathrm{Th}_{\Delta_1}}(T_1,T'_1)\quad\text{pl. fidèle.}
\]
L'image essentielle est formée des foncteurs $T_1-T'_1$ qui sont tels que
\[
S=T(\mathrm{Ens})\inserted{=T_1(\mathrm{Ens})}\longrightarrow S'=T'(\mathrm{Ens})=T'_1(\mathrm{Ens})
\]
transforme objets de P.F. en objets de P.F. [ On pourrait les appeler les morphismes de présent. finie, ou cohérents, de théories \ldots ].

% vertical emphasis stroke beside the question
\uline{Question} Considérons $\Delta$ définie par $d'=$ ens des diagrammes finis \uline{discrets}, $d''=\varnothing$, pas de condition d'exactitude. Soit $R$ une petite cat $\in\mathrm{Ob}\,\Delta$, $\Sigma=R^\circ\subset S=\unsure{T_R|\mathrm{Th}_R}(\mathrm{Ens})=\underline{\mathrm{Hom}}_{\mathrm{prod}}(R,\mathrm{Ens})$. Peut-il se faire que le $=$ $S$ soit défini par des sous-catégories \deleted{essentielles}\inserted{str. pleines} différentes ? \deleted{Oui, comme on voit avec} Comment trouver \inserted{toutes} les $\Sigma$ qui marchent, i.e. tels que $S\cong\rho_{\Delta,\Delta_1}(\Sigma)$ ? on aura évidem\textsuperscript{mt} $\Sigma\subset S_{pF}=\Sigma_0$ et $\unsure{\Sigma_{pF}|\Sigma_0|S_{pF}}\cong\rho_{\Delta,\Delta_0}(\Sigma)$ (condition néc. et suff.) Donc la question se repose en termes d'un $\Sigma_0\in\mathrm{Ob}\,\Delta_0$ et d'une \deleted{famille \unsure{d'obj.|génératrice} dans} \inserted{sous-catégorie pleine} $\Sigma$ de $\Sigma_0$ stable par sommes finies\ldots
% archive sheet number: 78
